\documentclass{beamer}

\usepackage{graphicx}
\usepackage{amsmath}

\title{Motivation for the Immersed Boundary Method}
\author{Laura Miller}
\date{\today}

\begin{document}

\frame{\titlepage}

\begin{frame}{Resources}
  \begin{itemize}
    \item \href{https://www.youtube.com/watch?v=PJyQA0vwbgU}{Resource 1: Introduction to IBM by Nick Battista}
    \item \href{https://drive.google.com/drive/folders/1G2bkEsv2Jz_R7165a5BqPbW4dU2hXBHb?usp=sharing}{Resource 2: Matea Santiago's lectures on IBM (mostly related to discretization)}
  \end{itemize}
\end{frame}

\begin{frame}{What is the Immersed Boundary Method?}
    \begin{itemize}
        \item The immersed boundary method is a computational approach for simulating fluid-structure interaction (FSI) problems.
        \item It was originally developed by Charlie Peskin in the 1970s to study the flow of blood in the heart.
        \item The method uses a Lagrangian description of the immersed structure (e.g., a heart valve) and an Eulerian description of the fluid (e.g., blood).
    \end{itemize}
\end{frame}

\begin{frame}{Applications of the Immersed Boundary Method}
\footnotesize
The immersed boundary (IB) method has been applied to a wide range of
fluid--structure interaction (FSI) problems where flexible or moving
structures interact with a viscous fluid.

\medskip
\textbf{Classical and biological applications:}
\begin{itemize}
    \item Blood flow in the heart and vasculature (valves, chambers, embryonic hearts)
    \item Swimming and propulsion of fish, jellyfish, and microorganisms
    \item Cilia- and flagella-driven flows
    \item Respiratory flows involving flexible airways
\end{itemize}

\medskip
\textbf{Engineering and physical systems:}
\begin{itemize}
    \item Flapping flags, membranes, and flexible beams
    \item Flow-induced vibrations of elastic structures
    \item Mixing and transport driven by moving boundaries
\end{itemize}

\medskip
\textbf{Emerging and interdisciplinary applications:}
\begin{itemize}
    \item Soft robotics and bio-inspired propulsion
    \item Tissue--fluid interactions and mechanobiology
    \item Porous and compliant surfaces in turbulent flows
\end{itemize}

\end{frame}

\begin{frame}{Advantages and Disadvantages of the IB Method}
\small

\textbf{Advantages:}
\begin{itemize}
    \item Handles complex, moving, and highly deformable geometries
    \item Avoids remeshing by using a fixed Cartesian fluid grid
    \item Naturally accommodates large structural deformations
    \item Unified framework for fibers, surfaces, and solids
    \item Well-suited for biological and soft-matter systems
\end{itemize}

\medskip
\textbf{Disadvantages and challenges:}
\begin{itemize}
    \item Interface is smeared over a few grid cells (limited sharpness)
    \item Time-step restrictions for stiff elastic forces
    \item Numerical slip at the fluid--structure interface
    \item Higher computational cost than rigid-boundary solvers
\end{itemize}

\medskip
\textbf{Important note:}
Many of these limitations have motivated modern extensions such as
implicit coupling, adaptive mesh refinement, and finite-element
immersed boundary methods.
\end{frame}

\begin{frame}{Key Idea of the Immersed Boundary Method}
\small

The immersed boundary method models fluid--structure interaction
by coupling a fluid solver with a force-based representation
of the immersed structure.

\medskip
\textbf{Core concepts:}
\begin{itemize}
    \item The fluid is solved on a fixed Eulerian grid
    \item The structure is represented by Lagrangian points or elements
    \item Structural forces are computed from elasticity or constitutive laws
    \item These forces are spread to the fluid as body forces
    \item The fluid velocity is interpolated back to move the structure
\end{itemize}

\medskip
\textbf{Physical interpretation:}
\begin{itemize}
    \item The structure influences the fluid through localized forces
    \item The fluid transports and deforms the structure
    \item No explicit boundary tracking is required
\end{itemize}

\medskip
\textbf{In essence:}
\[
\text{Fluid dynamics} \;\;+\;\; \text{elastic forces} \;\;\Rightarrow\;\; \text{coupled motion}
\]
\end{frame}

\begin{frame}{What the Immersed Boundary Method Is (and Is Not)}
\small

\textbf{What IBM \emph{is}:}
\begin{itemize}
    \item A \emph{force-based} formulation of fluid--structure interaction
    \item A method that couples Lagrangian structures to an Eulerian fluid
          through integral transforms
    \item A framework that enforces kinematic constraints \emph{weakly}
          via regularized delta functions
\end{itemize}

\medskip

\textbf{What IBM \emph{is not}:}
\begin{itemize}
    \item Not a sharp-interface or body-fitted method
    \item Not a pointwise enforcement of no-slip boundary conditions
    \item Not restricted to small deformations or simple geometries
    \item Not ignoring stresses or elasticity (these are encoded in the
          Lagrangian force density)
\end{itemize}

\medskip

\textbf{Key perspective:}  
IBM trades sharp geometric resolution for robustness, simplicity, and
the ability to handle large, complex motions.

\end{frame}

\begin{frame}{When Is the Immersed Boundary Method a Good Choice?}
\small

The immersed boundary method is particularly effective when:
\begin{itemize}
    \item The structure undergoes large deformations or topological changes
    \item The geometry is complex or time-dependent
    \item Exact enforcement of sharp boundaries is less critical than robustness
    \item The physics is dominated by viscous or moderately inertial flows
\end{itemize}

\medskip
It is less well-suited for:
\begin{itemize}
    \item Very high Reynolds number flows with thin boundary layers
    \item Problems requiring exact stress resolution at sharp interfaces
\end{itemize}
\end{frame}

\begin{frame}{Accuracy and Weak Enforcement of Boundary Conditions}
\small

In the immersed boundary method, boundary conditions are enforced
\emph{weakly} rather than pointwise.

\medskip

\textbf{Consequences:}
\begin{itemize}
    \item Velocity is continuous across the immersed structure
    \item Pressure and stress may be discontinuous
    \item No-slip is satisfied in an averaged (integral) sense
\end{itemize}

\medskip

\textbf{Accuracy depends on:}
\begin{itemize}
    \item Eulerian grid spacing $h$
    \item Lagrangian mesh spacing
    \item Smoothness and support of the regularized delta function
\end{itemize}

\medskip

As $h \rightarrow 0$ and the Lagrangian discretization is refined,
the regularized formulation converges to the continuum IB equations.

\medskip

\textbf{Interpretation:}  
IBM prioritizes stability and robustness over sharp interface resolution,
which is often appropriate for soft, moving, or biological structures.

\end{frame}

\begin{frame}{Common Failure Modes and Practical Pitfalls}
\small

While robust, the immersed boundary method has several common pitfalls.

\medskip

\textbf{Discretization issues:}
\begin{itemize}
    \item Lagrangian spacing too coarse $\Rightarrow$ force leakage
    \item Lagrangian spacing too fine $\Rightarrow$ stiff force response
    \item Mismatch between Eulerian and Lagrangian resolutions
\end{itemize}

\medskip

\textbf{Kernel-related issues:}
\begin{itemize}
    \item Using different kernels for spreading and interpolation
    \item Kernels that violate moment conditions
    \item Insufficient smoothness leading to grid-scale noise
\end{itemize}

\medskip

\textbf{Time-stepping issues:}
\begin{itemize}
    \item Explicit treatment of stiff elastic forces
    \item Severe time-step restrictions for strong elasticity
\end{itemize}

\medskip

\textbf{Takeaway:}  
Many IBM failures are not conceptual—they arise from inconsistent
discretization choices.

\end{frame}


\begin{frame}{Immersed Boundary Method vs.\ Other FSI Approaches}

\small
Fluid--structure interaction problems can be approached in several ways,
depending on geometry, deformation, and accuracy requirements.

\tiny
\medskip
\begin{center}
\begin{tabular}{l|c|c|c}
\textbf{Method} & \textbf{Geometry Handling} & \textbf{Deformation} & \textbf{Interface} \\
\hline
\textbf{Immersed Boundary (IB)} 
& Fixed Cartesian grid 
& Large, complex 
& Diffuse (regularized) \\
\textbf{ALE / Body-Fitted Mesh} 
& Moving, fitted mesh 
& Small--moderate 
& Sharp \\
\textbf{Immersed Interface / Ghost Fluid} 
& Fixed grid 
& Small--moderate 
& Sharp \\
\textbf{Fictitious Domain / Penalization} 
& Fixed grid 
& Moderate 
& Diffuse or penalized \\
\end{tabular}
\end{center}

\small
\medskip
\textbf{Key tradeoffs:}
\begin{itemize}
    \item IB prioritizes robustness and geometric flexibility over sharp interface resolution
    \item Body-fitted methods offer higher accuracy near interfaces but require remeshing
    \item Sharp-interface immersed methods improve accuracy but are harder to implement
\end{itemize}

\medskip
\textbf{Bottom line:} IB excels when geometry moves and deforms more than accuracy demands sharp boundaries.
\end{frame}

\begin{frame}{Grid-Fitted (Body-Fitted) FSI Methods: Accuracy vs.\ Cost}
\small

In grid-fitted (or body-fitted) FSI methods, the computational mesh
conforms exactly to the fluid--structure interface.

\medskip
\textbf{Key features:}
\begin{itemize}
    \item Fluid mesh aligns with the structure boundary
    \item Sharp enforcement of boundary conditions (no-slip, traction continuity)
    \item High accuracy near the interface
\end{itemize}

\medskip
\textbf{Typical approaches:}
\begin{itemize}
    \item Arbitrary Lagrangian--Eulerian (ALE) methods
    \item Moving mesh finite element / finite volume methods
    \item Coupled fluid--solid solvers with interface matching
\end{itemize}
\end{frame}

\begin{frame}{Grid-Fitted FSI Methods}
\textbf{Primary costs and challenges:}
\begin{itemize}
    \item Mesh deformation and distortion as the structure moves
    \item Frequent remeshing for large deformations
    \item Complex data transfer between fluid and solid solvers
    \item High implementation and maintenance complexity
\end{itemize}

\medskip
\textbf{Bottom line:}
\[
\text{Grid-fitted methods trade geometric flexibility for interface accuracy.}
\]

They are most effective when:
\begin{itemize}
    \item Deformations are small to moderate
    \item Interface accuracy is critical
    \item Geometry changes are limited or predictable
\end{itemize}

\end{frame}

\begin{frame}{The Immersed Interface Method (IIM)}
\footnotesize

The immersed interface method is a fixed-grid approach for
fluid--structure interaction and multiphase flows with \textbf{sharp interfaces}.

\medskip
\textbf{Core idea:}
\begin{itemize}
    \item Use a Cartesian Eulerian grid (no body-fitted mesh)
    \item Represent the interface explicitly (often via level sets or marker points)
    \item Modify finite difference stencils near the interface to enforce jump conditions
\end{itemize}

\medskip
\textbf{Key distinction from IBM:}
\begin{itemize}
    \item IBM spreads forces using regularized delta functions
    \item IIM enforces \emph{exact jump conditions} across the interface
\end{itemize}

\medskip
\textbf{Typical jump conditions enforced:}
\begin{itemize}
    \item Velocity continuity
    \item Discontinuous pressure
    \item Discontinuous stress proportional to interface forces
\end{itemize}

\medskip
\textbf{Result:} Sharp interface representation on a fixed grid.

\end{frame}

\begin{frame}{Immersed Interface Method: Strengths and Tradeoffs}
\footnotesize

\textbf{Advantages:}
\begin{itemize}
    \item Sharp resolution of interfaces (no numerical smearing)
    \item Higher accuracy near boundaries than regularized IB
    \item Exact enforcement of jump conditions
    \item Fixed Cartesian grid (no remeshing)
\end{itemize}

\medskip
\textbf{Costs and limitations:}
\begin{itemize}
    \item Complex implementation
    \item Requires analytic expressions for jump conditions
    \item Difficult for large deformations or topological changes
    \item Less flexible for highly elastic or active structures
\end{itemize}

\medskip
\textbf{Best suited for:}
\begin{itemize}
    \item Interfaces with known physics and jump conditions
    \item Problems requiring high accuracy near boundaries
    \item Moderate interface motion without severe distortion
\end{itemize}

\medskip
\textbf{Contrast with IBM:}
\[
\text{IIM prioritizes interface accuracy; IBM prioritizes robustness and flexibility.}
\]

\end{frame}

\begin{frame}{A Major Practical Advantage: Open-Source IBM Software}
\small

One of the strongest advantages of the immersed boundary method is the
availability of mature, open-source software—including for full 3D simulations.

\medskip

\textbf{Well-developed open-source IBM frameworks include:}
\begin{itemize}
    \item \textbf{IBAMR}: large-scale, parallel 3D simulations with AMR,
          finite differences, and finite elements
    \item \textbf{IB2d / IB2d-reloaded}: lightweight, accessible 2D implementations
          for research and teaching
    \item Research codes released alongside published studies
\end{itemize}

\medskip

\textbf{Why this matters:}
\begin{itemize}
    \item Low barrier to entry for new researchers and students
    \item Full access to algorithms, discretizations, and solver details
    \item Reproducibility and transparency in computational studies
    \item Rapid prototyping of new models and coupling strategies
\end{itemize}

\end{frame}

\begin{frame}{Open Source FSI}

\textbf{Contrast with many FSI methods:}
\begin{itemize}
    \item Grid-fitted and ALE methods often rely on proprietary solvers
    \item Sharp-interface methods may lack robust, public 3D implementations
    \item Significant development effort required before doing science
\end{itemize}

\medskip

\textbf{Takeaway:}  
IBM is not only a mathematical framework—it is a well-supported
\emph{computational ecosystem} for 3D fluid--structure interaction.

\end{frame}


\end{document}